\chapter{Timing}
\label{ch:timing}

\section{Clocks}

The device runs from the system's master clock. It does one PPU cycle every two master clocks, and one pixel every two
PPU cycles, which is four master clocks. Table~\ref{tab:clocks} gives the rates for the master clock of the SIRCUS
system, 21,477,272~Hz, which is the NTSC colour subcarrier frequency multiplied by six.

\begin{table}[H]
	\centering
	\begin{tabular}{lrr}
		\toprule
		\textbf{Clock} & \textbf{Master clocks} & \textbf{At 21,477,272 Hz} \\
		\midrule
		Master clock   & 1                      & 21,477,272 Hz             \\
		PPU cycle      & 2                      & 10,738,636 Hz             \\
		Pixel          & 4                      & 5,369,318 Hz              \\
		Line           & 1364                   & 15,746 Hz                 \\
		Frame          & 357,368                & 60.1 Hz                   \\
		\bottomrule
	\end{tabular}
	\caption{Clock Rates}
	\label{tab:clocks}
\end{table}

\section{The CPU Clock}
\label{sec:cpu-clock}

The CPU advances one execution phase on each rising edge of \texttt{CLKI} (\hwmanual{}). The system supplies the CPU
and the video device with the same \texttt{CLKI}, which is the only clock the device has, and \texttt{CLKI} is one PPU
cycle: the master clock divided by two, 10,738,636~Hz at 21,477,272~Hz. The figures that follow are consequences:

\begin{itemize}
	\item An ordinary instruction takes 6 \texttt{CLKI} cycles, 12 master clocks, so the CPU can run about 1.79 million
	      instructions a second.
	\item A line is 682 \texttt{CLKI} cycles, about 113 instructions, and a frame is 178,684 cycles. NTSC vertical blanking, 38
	      lines, is 25,916 cycles, about 4,300 instructions.
	\item The registers are sampled 228 master clocks, 114 cycles, after the start of a line. A line-compare interrupt is
	      recognised at an instruction boundary, so the handler's first instruction is fetched at best 6 and at worst about 12
	      cycles later (\archmanual{}, interrupt latency; more after a coprocessor call or a DMA transfer). That leaves about 100
	      cycles, some 16 instructions, for a handler to change a register before it takes effect on the same line.
	\item A \texttt{DMAT} of 255 words takes 522 cycles, three quarters of a line, and cannot be interrupted.
\end{itemize}

The video device never waits for the CPU, and it gives the CPU no wait states: every access to the device completes in
the cycle it is made (Section~\ref{sec:memory-vram-access}).

\section{Line Timing}

A line is 1364 master clocks, 682 PPU cycles, in three parts. A line begins with the horizontal sync pulse
(Section~\ref{sec:video-output}), and the picture follows the back porch; the front porch ends the line. The line
timing is fixed and the same in both modes:

\begin{table}[H]
	\centering
	\begin{tabular}{lrrl}
		\toprule
		\textbf{Part}       & \textbf{Master clocks} & \textbf{PPU cycles} & \textbf{Purpose}                    \\
		\midrule
		Sync and back porch & 260                    & 130                 & Nothing is drawn; Mode 2 prefetches \\
		Visible             & 1024                   & 512                 & 256 pixels at 4 master clocks each  \\
		Front porch         & 80                     & 40                  & Nothing is drawn                    \\
		\midrule
		\textbf{Line}       & \textbf{1364}          & \textbf{682}        &                                     \\
		\bottomrule
	\end{tabular}
	\caption{Line Timing}
	\label{tab:line-timing}
\end{table}

The visible time of a line is a whole number of master clocks per pixel: four.

\begin{figure}[H]
	\centering
	\begin{tikzpicture}[
			x=0.0057\textwidth,
			y=0.62cm,
			bar/.style={draw, minimum height=0.55cm, font=\scriptsize, align=center, inner sep=1pt},
			lane/.style={font=\scriptsize, anchor=east},
			marker/.style={dashed, thick, sircgray},
			note/.style={font=\scriptsize, text=sircgray}
		]
		% Line parts (1364 master clocks)
		\fill[sirclightgray] (0,0) rectangle (26,0.55);
		\fill[white] (26,0) rectangle (128.4,0.55);
		\fill[sirclightgray] (128.4,0) rectangle (136.4,0.55);
		\draw (0,0) rectangle (136.4,0.55);
		\draw (26,0) -- (26,0.55);
		\draw (128.4,0) -- (128.4,0.55);
		\node[font=\scriptsize] at (13,0.28) {sync, back porch};
		\node[font=\scriptsize] at (77.2,0.28) {visible, 1024 clocks};
		\node[font=\tiny] at (132.4,0.28) {front};
		% Lanes
		\node[lane] at (-1,1.4) {Background fetch};
		\draw[fill=blue!15] (22.8,1.1) rectangle (26,1.7);
		\draw[fill=blue!8] (26,1.1) rectangle (128.4,1.7);
		\node[font=\scriptsize] at (77.2,1.4) {fetch a group of four pixels every 16 master clocks};
		\node[lane] at (-1,2.5) {Sprite fetch};
		\draw[fill=orange!20] (0,2.2) rectangle (22.8,2.8);
		\draw[fill=orange!20] (128.4,2.2) rectangle (136.4,2.8);
		\node[font=\scriptsize] at (11.4,2.5) {fetch sprite rows};
		\node[lane] at (-1,3.6) {Sprite scan};
		\draw[fill=green!10] (26,3.3) rectangle (128.4,3.9);
		\node[font=\scriptsize] at (77.2,3.6) {scan the sprite table for the next line (internal memory)};
		% Markers
		\draw[marker] (0,-0.2) -- (0,4.6);
		\draw[marker] (22.8,-0.2) -- (22.8,4.6);
		\draw[marker] (26,-0.2) -- (26,4.6);
		\node[note, anchor=south west] at (0,4.65) {line-compare IRQ4};
		\node[note, anchor=south] at (22.8,5.25) {registers sampled};
		\node[note, anchor=south west] at (26.6,4.65) {first pixel};
		% Axis
		\foreach \x/\l in {0/0, 26/260, 128.4/1284, 136.4/1364} { \node[font=\tiny, anchor=north] at (\x,-0.2) {\l}; }
	\end{tikzpicture}
	\caption{Timeline of One Visible Line, in Master Clocks}
	\label{fig:line-timeline}
\end{figure}

\section{Frame Timing}

A frame is 262 lines on NTSC and 312 on PAL. Lines 0 to 223 are visible and draw the picture. The remaining lines, 38
on NTSC and 88 on PAL, are the vertical blanking interval, in which nothing is drawn and the CPU may update video
memory (Section~\ref{sec:memory-vram-access}). The device fetches nothing in these lines.

Table~\ref{tab:times} gives the times. They follow from the clocks: on NTSC the PPU cycle is 93.1~ns and a pixel, four
master clocks, is 186.2~ns.

\begin{table}[H]
	\centering
	\begin{tabular}{lrr}
		\toprule
		                       & \textbf{NTSC} & \textbf{PAL} \\
		\midrule
		Master clock           & 21.477~MHz    & 21.281~MHz   \\
		Line                   & 63.5~$\mu$s   & 64.1~$\mu$s  \\
		Frame                  & 16.64~ms      & 20.00~ms     \\
		Frame rate             & 60.10~Hz      & 50.00~Hz     \\
		Vertical blanking      & 2.41~ms       & 5.64~ms      \\
		Pixel                  & 186.2~ns      & 188.0~ns     \\
		Visible part of a line & 47.7~$\mu$s   & 48.1~$\mu$s  \\
		\bottomrule
	\end{tabular}
	\caption{Video Timing}
	\label{tab:times}
\end{table}

\subsection{PAL}

The \texttt{OUT} bit of the Status register says which the device is. PAL has the same picture, 256 by 224 pixels in
lines 0 to 223, at the top, and the same line timing of 1364 master clocks, but 312 lines per frame, so the vertical
blanking interval is longer. A program written for NTSC runs unchanged on PAL, with more time in vertical blanking and
a lower frame rate. With a PAL master clock of 21,281,370 Hz (the PAL colour subcarrier multiplied by 4.8) a frame is
425,568 master clocks, 50.0 Hz. The region is fixed when the device is made and is read from \texttt{OUT}.

\section{Interrupts}
\label{sec:interrupts}

The device has two interrupt sources, on different lines, so a handler knows which it is serving. The device races the
beam: it never waits for the CPU, so an interrupt is raised at a fixed point of the frame whether or not the CPU is
ready, and the time the CPU takes to respond is the program's concern. A program that shares \activelow{IRQ4} with
another device can read the \texttt{LIP} bit of the Status register to see whether the device's interrupt is the one
pending.

\subsection{Vertical Sync}

The vertical sync interrupt is the non-maskable interrupt, \activelow{NMI}, Level 5. It is edge-triggered, so the
device needs no acknowledgement: when \texttt{VIE} is set in the Interrupt Control register
(Section~\ref{sec:reg-intctl}), the device asserts the \activelow{NMI} pin for one PPU cycle at the first cycle of line
224, the start of vertical blanking. The program has the whole of vertical blanking to update video memory
(Section~\ref{sec:memory-vram-access}). The pulse is exactly one \texttt{CLKI} cycle (Section~\ref{sec:cpu-clock}). The
CPU samples \activelow{NMI} on every rising edge of \texttt{CLKI} except during reset (\hwmanual{}), so it does not
miss the pulse. A pulse that arrives while the handler of the previous one is still running, or while a Level 5 request
is pending, is not taken as an interrupt: it raises the Level Five Hardware Exception Conflict fault (\archmanual{}).
The handler must therefore return before the next vertical sync, 16.6~ms later on NTSC and 20.0~ms on PAL.

The \activelow{NMI} cannot be masked, so a program must install its handler before it sets \texttt{VIE}. A DMA transfer
cannot be interrupted, the \activelow{NMI} included. The edge is latched and the interrupt is taken when the transfer
ends, so it is delayed, not lost, but a transfer that spans line 224 shortens the time the program has in vertical
blanking by the delay. A program keeps its transfers short (Section~\ref{sec:cpu-clock} gives their lengths), or
arranges for them to finish before line 224.

\subsection{Line Compare}
\label{sec:line-interrupt}

The line-compare interrupt is interrupt level 4, \activelow{IRQ4}, the highest of the maskable levels. When
\texttt{LIE} is set, the device asserts the line at the first cycle of the line whose number is in the \texttt{LINE}
register (Section~\ref{sec:reg-line}). It fires on any line of the frame, visible or not, so lines 0 to 223 can be used
for raster effects. The \activelow{IRQ4} line is level-sensitive, so the device holds it asserted until the program
acknowledges it by writing any value to \texttt{LINE}, which also sets the line for the next interrupt, or by clearing
\texttt{LIE}. The \texttt{LIP} bit of the Status register is set while the interrupt is pending. A \texttt{LINE} value
at or beyond the number of lines in the frame never matches. If the line is reached again while an interrupt is still
pending, nothing is queued: the line is already asserted.

The palette is writable at any time and a write changes the pixels drawn after it; the registers are sampled once per
line, 32 master clocks before the first visible pixel in Mode 2 and Mode 3, so a register written within the 228 master
clocks after the interrupt takes effect on that line, and later ones on the next. The device does not wait for the CPU,
so how soon after the interrupt the CPU's write arrives is the program's concern, and a handler for a raster effect
must be short and of known length.

\section{Fetch Budget}

Reading video memory is the device's limit: one read per PPU cycle, eight per group of four pixels. The modes spend the
eight slots differently.

\begin{table}[H]
	\centering
	\begin{tabular}{lccc}
		\toprule
		\textbf{Use of the slots}                             & \textbf{Mode 1} & \textbf{Mode 2}    & \textbf{Mode 3}   \\
		\midrule
		Background tilemap entries                            & 4               & 2--4               & 3--4              \\
		Background tile data words                            & 4               & 2--4               & 3--4              \\
		Sprites                                               & 0               & 0                  & 0                 \\
		\midrule
		\textbf{Total, typical (Mode 2, 3) or fixed (Mode 1)} & \textbf{8}      & \textbf{4.75 of 8} & \textbf{6.4 of 8} \\
		\textbf{Total, worst case}                            & \textbf{8}      & \textbf{8}         & \textbf{8}        \\
		\bottomrule
	\end{tabular}
	\caption{Fetch Slots per Group of Four Pixels}
	\label{tab:fetch-budget}
\end{table}

\begin{figure}[H]
	\centering
	\begin{tikzpicture}[
			cell/.style={draw, minimum width=1.7cm, minimum height=0.85cm, font=\scriptsize, align=center, inner sep=1pt},
			flatc/.style={cell, fill=orange!20},
			isoc/.style={cell, fill=blue!15},
			isob/.style={cell, fill=blue!15, dashed},
			free/.style={cell, fill=gray!8, dashed},
			rowlab/.style={font=\small, anchor=east},
			head/.style={font=\scriptsize, text=sircgray}
		]
		\foreach \i in {0,...,7} { \node[head] at (\i*1.75,0.7) {slot \i}; }
		\node[rowlab] at (-1.0,0) {Mode 1};
		\foreach \i/\t in {0/B1 entry,1/B2 entry,2/B3 entry,3/B4 entry,4/B1 word,5/B2 word,6/B3 word,7/B4 word} { \node[flatc] at (\i*1.75,0) {\t}; }
		\node[rowlab] at (-1.0,-1.1) {Mode 2};
		\foreach \i/\t in {0/B1 entry A,1/B1 word A,4/B2 entry A,5/B2 word A} { \node[isoc] at (\i*1.75,-1.1) {\t}; }
		\foreach \i/\t in {2/B1 entry B,3/B1 word B,6/B2 entry B,7/B2 word B} { \node[isob] at (\i*1.75,-1.1) {\t}; }
		\node[rowlab] at (-1.0,-2.2) {Mode 3};
		\foreach \i/\t in {0/B1 entry,1/B1 word,2/B2 entry,3/B2 word} { \node[flatc] at (\i*1.75,-2.2) {\t}; }
		\foreach \i/\t in {4/B3 entry A,5/B3 word A} { \node[isoc] at (\i*1.75,-2.2) {\t}; }
		\foreach \i/\t in {6/B3 entry B,7/B3 word B} { \node[isob] at (\i*1.75,-2.2) {\t}; }
		% Legend
		\node[flatc, minimum width=1.4cm] at (0,-3.4) {orthogonal};
		\node[isoc, minimum width=1.4cm] at (1.75,-3.4) {isometric};
		\node[isob, minimum width=1.4cm] at (3.5,-3.4) {second cell};
		\node[head, anchor=west] at (5.0,-3.4) {dashed: used only if needed};
	\end{tikzpicture}
	\caption{Who Uses Each Fetch Slot of a Group of Four Pixels}
	\label{fig:slot-map}
\end{figure}

The status layer of Mode 2 takes its reads from the slots the isometric backgrounds leave spare, through a queue
(Section~\ref{sec:m2-status-layer}); they are not in the table. Mode 3's orthogonal backgrounds read a tilemap entry
and a tile word in every group, at fixed slots (Chapter~\ref{ch:mode3}); with 2-bit tiles they read only in every
second group. A mode with more backgrounds or a different mix of fetches than these needs more reads per group, such a
queue, or a faster memory.

\section{Pixel Pipeline}

The device draws a group of four pixels in two stages. During one group the device fetches the data for a later group
and, at the start of the group, combines the data it already has into four finished pixels, which are shown one at a
time over the group. Mode 1 combines the fetches of the previous group; Mode 2 combines the data it prefetched
(Section~\ref{sec:m2-fetch}). The first pixels of a line in Mode 1 are at the position the scroll registers give
(Section~\ref{sec:m1-scrolling}).

Figure~\ref{fig:fetch-timing} shows the schedule of one group in Modes 1 and 3, one read to a \texttt{CLKI} cycle, with
the four pixels that are shown while it is made. A pixel lasts two cycles. In Mode 3 the reads of the second cell of
the isometric background are made only when the group straddles two cells.

\begin{figure}[H]
	\centering
	\small
	\begin{tikztimingtable}[sirc timing diagram]
		CLKI                  & L 7{0.5H 0.5L}                                                  \\
		Slot                  & D{0} D{1} D{2} D{3} D{4} D{5} D{6} D{7}                         \\
		Mode 1 read           & D{B1e} D{B2e} D{B3e} D{B4e} D{B1w} D{B2w} D{B3w} D{B4w}         \\
		Mode 3 read           & D{B1e} D{B1w} D{B2e} D{B2w} D{B3e} D{B3w} D{B3e'} D{B3w'}       \\
		Pixel shown           & 2D{p0} 2D{p1} 2D{p2} 2D{p3}                                     \\
	\end{tikztimingtable}
	\caption{Fetch Schedule of One Group of Four Pixels (e: Tilemap Entry, w: Tile Word, $'$: Second Cell)}
	\label{fig:fetch-timing}
\end{figure}
